\subsection{Primitive reflection singularities for arbitrary mass}
\label{sec:primitive-time}

At zero momentum, the full gamma-function envelope in
\eqref{eq:detuned-spectrum} can be
retained when locating a particular family of complex-time singularities.
Define \(s_\pm=(1\pm\ii)\omega/4\) and
\begin{equation}
 B_\nu(\omega)=8^\nu
 \frac{\Gamma(h+s_+)\Gamma(h+s_-)}
      {\Gamma(1-h+s_+)\Gamma(1-h+s_-)} .
 \label{eq:analytic-envelope}
\end{equation}
It equals \(\omega^{2\nu}P_\nu(\omega/2)\) on positive real \(\omega\).
The numerator gamma functions have no poles in the closed sector
\(|\arg\omega|\leq\pi/2\), and the reciprocal gamma functions are entire
in their arguments. In this sector Stirling asymptotics give
\(B_\nu(\omega)/\omega^{2\nu}\to1\) uniformly at infinity, with the
power continued from real positive \(\omega\).

Let \(\tau=\beta/2\), \(\varphi_\nu=\pi(\nu+1)\), and
\(\rho_{\rm vac}(\omega)=C_\nu\omega^{2\nu}\) for \(\omega>0\).
The exact thermal Wightman function has a high-frequency sum over
the two geometric factors. With any fixed cutoff \(\omega_0>0\),
introduce
\begin{equation}
 \Phi_\nu(L)=\int_{\omega_0}^\infty
      \dd\omega\, e^{-L\omega}B_\nu(\omega).
\end{equation}
This converges initially for \({\rm Re}\,L>0\). For
\({\rm Im}\,L>0\), rotate the frequency tail to the negative
imaginary axis and include the finite arc from \(\omega_0\) to
\(-\ii\omega_0\). For \({\rm Im}\,L<0\), rotate to the positive
imaginary axis. No poles cross either deformation. Thus
\(\Phi_\nu\) continues on the \(L\) plane slit along the nonpositive
real axis. Its leading singular behaviour on this branch is
\begin{equation}
 \Phi_\nu(L)\sim\frac{\Gamma(2\nu+1)}{L^{\,2\nu+1}},
 \qquad L\longrightarrow0
 \label{eq:laplace-asymptotic}
\end{equation}
on closed subsectors of the slit plane. Rescaling the integral
on a convergent ray gives \eqref{eq:laplace-asymptotic}. The finite
arc and any change of the lower cutoff are analytic in \(L\).

On the thermal strip, the part of the two-sided frequency integral
with \(|\omega|\geq\omega_0\) is
\begin{align}
 G^>_{\rm tail}(t)
 &=\frac{C_\nu}{2\pi}\sum_{j,k\geq0}
 e^{-\ii(j-k)\varphi_\nu}\!
 \left[\Phi_\nu(L^+_{jk})+\Phi_\nu(L^-_{jk})\right],\notag\\
 L^+_{jk}&=(j+k)\tau+\ii[t-(j-k)\tau],\notag\\
 L^-_{jk}&=(j+k+2)\tau-\ii[t+(j-k)\tau].
 \label{eq:laplace-image}
\end{align}
The finite-frequency remainder is entire in \(t\), since it
integrates an integrable spectral kernel over a bounded interval.
The sum over sufficiently large \(j+k\) converges locally uniformly:
the cutoff supplies exponential suppression as \({\rm Re}\,L^\pm\)
increases, while the number of terms per \(j+k\) grows only linearly.
The continuation therefore reduces to a finite number of rotated
Laplace integrals in any bounded region of complex time.

Define the primitive reflection times
\begin{equation}
 t_n^\pm=\pm n\tau+\ii n\tau,\qquad n=1,2,\ldots .
\end{equation}
These are the first upper singularities along their respective
vertical lines. Approach each from the thermal strip at fixed
\({\rm Re}\,t=\pm n\tau\). For a positive-frequency term with
\(j-k\ne\pm n\), \(L^+_{jk}\) has nonzero imaginary part along
the path, so the continuation avoids its slit. When \(j-k=\pm n\),
one has \(j+k\geq n\), and exactly one term has \(L^+_{jk}=0\) at
the target. Every negative-frequency \(L^-_{jk}\) has positive
real part near the target. Hence the full continued correlator
has the leading contribution
\begin{equation}
 G^>(t)\sim
 \frac{C_\nu\,\Gamma(2\nu+1)}{2\pi}\,
 \frac{e^{\mp\ii n\varphi_\nu}}
      {[\ii(t-t_n^\pm)]^{2\nu+1}} .
 \label{eq:primitive-coefficient}
\end{equation}
For nonintegral \(2\nu+1\) this is a branch point. At integral
\(2\nu+1\) it gives the leading pole power; weaker terms can
depend on \(\nu\). The result identifies the primitive echoes of
the full correlator, including the \(n=2\) pair, without splitting
the envelope into perturbative and exponentially small pieces.
Mixed images with \(j,k>0\) can lie on a cut from an earlier
primitive echo. Section~\ref{sec:mixed-time} continues the correlator
on the two sides of such cuts.

\begin{table}[t]
 \centering\small
 \begin{tabular}{ccccc}
 \hline
 \(n\) & \(\nu\) & \(\epsilon=0.16\) & \(\epsilon=0.08\) &
 \(\epsilon=0.04\) \\
 \hline
 1 & 0.75 & 2.834e-04 & 2.336e-05 & 3.662e-06 \\
2 & 0.75 & 2.848e-04 & 2.364e-05 & 3.706e-06 \\
2 & 1 & 1.560e-04 & 9.291e-06 & 9.676e-07 \\
2 & 1.25 & 1.155e-04 & 5.962e-06 & 2.859e-07 \\
3 & 1.25 & 1.155e-04 & 5.962e-06 & 2.859e-07 \\
\hline

 \end{tabular}
 \caption{Relative error in extracting the full leading coefficient
 in \eqref{eq:primitive-coefficient} at \(t=t_n^+-\ii\epsilon\).
 The integral keeps the exact gamma envelope, the finite-frequency
 remainder and both frequency signs.}
 \label{tab:time-check}
\end{table}

Table~\ref{tab:time-check} shows the numerical approach to the
coefficient. We rotate the necessary finite set of Laplace integrals
and sum the higher harmonics geometrically before integrating the
convergent tails. At the final \(\epsilon\), rotating by \(\pi/3\)
instead of \(\pi/2\) changes the scaled correlator by less than
\(1.3\times10^{-11}\) in these examples. The coefficient at
\(t_n^-\) is the complex conjugate of that at \(t_n^+\), as
expected from the real frequency spectral data.
